% ============================================================
% section5_theory.tex - Раздел: Теория проволочных поляризаторов
% ============================================================

\section{Теоретическая модель проволочного поляризатора}

В предыдущих разделах мы научились загружать, визуализировать и предварительно обрабатывать экспериментальные данные ТГц-спектроскопии. Однако для глубокого анализа угловой зависимости пропускания нам необходима теоретическая модель. В этом разделе мы познакомимся с физической и математической основой работы металлических проволочных решеток в ТГц-диапазоне, опираясь на классические статьи A.~Blanco (1986) и K.~Kaltenecker (2019).

\subsection{Физический смысл проволочной решетки}

Идеальный поляризатор пропускает 100\% излучения с поляризацией, перпендикулярной проволокам, и полностью блокирует излучение с поляризацией, параллельной им. В реальности проволочные решетки ограничены геометрическими размерами — шагом решетки $p$ и диаметром проволоки $d$. 

Когда длина волны падающего излучения $\lambda$ становится сопоставимой с шагом решетки $p$ (так называемая «дифракционная область», $p > \lambda$), эффективность поляризатора падает: появляется паразитное пропускание параллельной поляризации (обозначим его $t_\parallel$), а пропускание перпендикулярной ($t_\perp$) становится меньше единицы. Кроме того, возникают дифракционные эффекты и потери энергии на высшие моды, что требует более сложных математических моделей.

\subsection{Метод эквивалентных длинных линий (Transmission Line Analogy)}

Для описания взаимодействия электромагнитных волн с проволочной решеткой применяется мощный инженерный подход — аналогия с длинными линиями, предложенная Маркувицем (Marcuvitz). 
Согласно этой модели, электромагнитная волна, распространяющаяся в свободном пространстве, рассматривается как сигнал в волноводе, а проволочная решетка — как препятствие с определенным импедансом (комплексным сопротивлением).

\begin{itemize}
    \item \textbf{Емкостное препятствие (Capacitive obstacle):} Возникает, когда вектор электрического поля \textbf{перпендикулярен} проволокам.
    \item \textbf{Индуктивное препятствие (Inductive obstacle):} Возникает, когда вектор электрического поля \textbf{параллелен} проволокам.
\end{itemize}

Метод Бланко расширяет эту теорию на дифракционную область, вводя комплексные импедансы с ненулевой действительной частью. Это позволяет учесть диссипацию энергии из-за связи с затухающими (эванесцентными) волнами высших порядков.

\subsection{Расчет коэффициентов пропускания (Blanco, 1986)}

Модель позволяет вычислить амплитудные коэффициенты пропускания $T_\perp$ и $T_\parallel$ через набор эквивалентных импедансов:

\begin{align}
    T_\perp &= \frac{2 Z_1 Z_2 / Z_0^2}{(1 + Z_1/Z_0)(Z_b/Z_0 + Z_2/Z_0)} \label{eq:T_perp} \\
    T_\parallel &= \frac{2 Z_3 Z_4 / Z_0^2}{(1 + Z_3/Z_0)(Z_d/Z_0 + Z_4/Z_0)(1 + Z_d/Z_0)} \label{eq:T_par}
\end{align}

Здесь $Z_0$ — волновое сопротивление свободного пространства. Импедансы $Z_1, Z_2, Z_3, Z_4$, а также базовые компоненты $Z_a, Z_b, Z_c, Z_d$ являются функциями от шага решетки $p$, диаметра $d$ и длины волны $\lambda$. В статье Blanco приведены явные аналитические формулы для этих компонент в виде бесконечных рядов, которые на практике можно обрывать на 10-м члене с сохранением высокой точности. Энергетические коэффициенты пропускания вычисляются как квадраты модулей: $K_1 = |T_\perp|^2$ и $K_2 = |T_\parallel|^2$.

\subsection{Матричный формализм Джонса и система двух поляризаторов}

В реальном эксперименте с THz-TDS (как описано в статье Kaltenecker, 2019) часто используется система из \textbf{двух} поляризаторов для контролируемого ослабления пучка. Первый поляризатор можно вращать на угол $\theta$, а второй жестко зафиксирован и пропускает только одну поляризацию (например, вдоль оси $x$).

\smartfigure{fig0.1.png}{Принципиальная схема эксперимента с двумя проволочными поляризаторами на пути ТГц излучения.}{fig:5_3}

Прохождение света через такую оптическую систему элегантно описывается матрицами Джонса. Поляризатор представляется диагональной матрицей $\mathbf{P}$, а его поворот — матрицей поворота $\mathbf{R}(\theta)$:
\begin{equation}
\mathbf{P} = \begin{pmatrix} t_\perp & 0 \\ 0 & t_\parallel \end{pmatrix}, \quad \mathbf{R}(\theta) = \begin{pmatrix} \cos\theta & -\sin\theta \\ \sin\theta & \cos\theta \end{pmatrix}
\end{equation}

\note{В статье Kaltenecker амплитудные коэффициенты обозначены как $t_\perp$ и $t_\parallel$, что полностью соответствует $T_\perp$ и $T_\parallel$ в модели Бланко.}

Тогда электрическое поле на выходе из системы $\mathbf{E}_{out}$ связано со входным полем $\mathbf{E}_{in}$ (поляризованным по оси $x$) следующим образом:
\begin{equation}
\mathbf{E}_{out} = \mathbf{P} \mathbf{R}(\theta) \mathbf{P} \mathbf{R}(-\theta) \mathbf{E}_{in}
\end{equation}
Здесь $\mathbf{R}(-\theta)$ переводит вектор поля в систему координат повернутого первого поляризатора, первая матрица $\mathbf{P}$ применяет пропускание, $\mathbf{R}(\theta)$ возвращает вектор в лабораторную систему координат, а вторая $\mathbf{P}$ описывает пропускание второго (неподвижного) поляризатора.

В результате аналитического перемножения этих матриц мы получаем выражение для выходного поля:
\begin{equation}
E_{out} = E_{in} \begin{pmatrix} t_\perp(t_\perp \cos^2\theta + t_\parallel \sin^2\theta) \\ t_\parallel(t_\perp - t_\parallel)\sin\theta\cos\theta \end{pmatrix}
\end{equation}

\subsection{Зачем нам эти формулы?}
В следующем разделе мы переведем эти громоздкие формулы в элегантный Python-код. Модель Бланко позволит нам:
\begin{enumerate}
    \item Вычислять $T_\perp(\lambda)$ и $T_\parallel(\lambda)$ для любых параметров решетки $p$ и $d$.
    \item Описывать деполяризацию пучка и возникновение эллиптичности.
    \item Строить теоретические кривые пропускания мощности $P(\theta, \lambda)$ и напрямую сравнивать их с нашими экспериментальными данными из датасета.
\end{enumerate}

Таким образом, мы сможем не просто наблюдать графики, а проверять, насколько хорошо реальный терагерцовый поляризатор описывается классической электродинамической моделью эквивалентных длинных линий.